% NB10c -- Diagramas de cuerpo libre del biela-manivela del motor diesel.
% Reemplaza las Figuras 8-12 (DCL a mano) del informe de Alarcon & Chilatra (2025), reducidas a 3 cuerpos.
% Convencion: F_X = fuerza que el cuerpo "aguas arriba" ejerce sobre el "aguas abajo" en el pasador X.
\documentclass[tikz,border=4mm]{standalone}
\usepackage{amsmath}
\usetikzlibrary{arrows.meta,calc,patterns}
\begin{document}
% Estilo comun de los diagramas TikZ del modulo VCR (NB10b) -- paleta del proyecto
\definecolor{manivela}{HTML}{2A78D6}
\definecolor{inferior}{HTML}{4A3AA7}
\definecolor{control}{HTML}{1BAF7A}
\definecolor{superior}{HTML}{EB6834}
\definecolor{tinta}{HTML}{3A3A38}
\definecolor{metal}{HTML}{C3C2B7}
\definecolor{perdida}{HTML}{E34948}
\definecolor{fuerza}{HTML}{C0392B}
\tikzset{
  barra/.style={line width=3.2pt, line cap=round},
  pin/.style={circle, fill=white, draw=tinta, line width=0.9pt, inner sep=1.6pt},
  fijo/.style={circle, fill=tinta, inner sep=2.6pt},
  vec/.style={-{Stealth[length=2.6mm]}, line width=0.9pt},
  F/.style={-{Stealth[length=3mm]}, line width=1.2pt, fuerza},
  cota/.style={{Stealth[length=1.8mm]}-{Stealth[length=1.8mm]}, line width=0.5pt, tinta},
  lbl/.style={font=\small, text=tinta},
}
\newcommand{\tierra}[1]{\fill[pattern=north east lines, pattern color=tinta] ($(#1)+(-5,-3)$) rectangle ($(#1)+(5,-6)$);
  \draw[tinta, line width=0.8pt] ($(#1)+(-5,-3)$) -- ($(#1)+(5,-3)$);}

\begin{tikzpicture}[x=1cm,y=1cm]
  % (1) manivela
  \begin{scope}[shift={(0,0)}]
    \draw[barra, manivela] (0,0) -- (-1.3,0.9); \node[fijo] at (0,0) {}; \node[pin] at (-1.3,0.9) {};
    \node[lbl, below left] at (-0.05,-0.05) {$O$}; \node[lbl, left=3pt] at (-1.3,0.9) {$C$};
    \draw[F] (0,0) -- ++(0.7,-0.9) node[lbl, below, fuerza] {$\vec F_O$};
    \draw[F] (-1.3,0.9) -- ++(0.3,-1.0) node[lbl, below, fuerza] {$-\vec F_C$};
    \draw[-{Stealth}, line width=1pt, manivela] (0.5,0.1) arc (10:200:0.5); \node[lbl, manivela] at (0.75,0.65) {$T_O$};
    \node[lbl, font=\bfseries] at (-0.4,-1.8) {(1) cig\"ue\~nal $OC$};
    \node[lbl, font=\footnotesize] at (-0.4,-2.2) {cm en $O$ (contrapesado), $I_{cig}$};
  \end{scope}
  % (2) biela
  \begin{scope}[shift={(4.2,-0.6)}]
    \draw[barra, superior] (0,0) -- (0.7,3.0); \node[pin] at (0,0) {}; \node[pin] at (0.7,3.0) {};
    \node[lbl, left=3pt] at (0,0) {$C$}; \node[lbl, left=3pt] at (0.7,3.0) {$D$};
    \fill[superior] (0.21,0.9) circle (2.2pt); \node[lbl, right, font=\footnotesize] at (0.25,0.85) {cm ($l_{cm}$ desde $C$)};
    \draw[F] (0,0) -- ++(-0.3,1.0) node[lbl, left, fuerza] {$\vec F_C$};
    \draw[F] (0.7,3.0) -- ++(-0.3,-1.0) node[lbl, right, fuerza, pos=0.8] {$-\vec F_D$};
    \draw[F, tinta] (0.21,0.9) -- ++(0,-0.7) node[lbl, below right] {$m_b\vec g$};
    \draw[-{Stealth}, line width=0.9pt, superior] (0.75,1.2) arc (-60:160:0.3); \node[lbl, superior, right] at (0.95,1.45) {$I_b\ddot\beta$};
    \node[lbl, font=\bfseries] at (0.4,-1.2) {(2) biela $CD$};
  \end{scope}
  % (3) piston
  \begin{scope}[shift={(8.6,0.6)}]
    \fill[metal] (-1.0,-0.35) rectangle (1.0,0.55); \draw[tinta] (-1.0,-0.35) rectangle (1.0,0.55);
    \node[pin] at (0,0) {}; \node[lbl, below right] at (0,0) {$D$};
    \foreach \x in {-0.75,-0.25,0.25,0.75} \draw[F, perdida] (\x,1.3) -- (\x,0.6);
    \node[lbl, perdida] at (0,1.6) {$F_{gas}=(p-p_{c\acute arter})A_p$};
    \draw[F] (0,0) -- ++(0.3,1.0) node[lbl, right, fuerza, pos=0.55] {$\vec F_D$};
    \draw[F] (-1.9,0.1) -- (-1.05,0.1) node[lbl, above, fuerza, pos=0.3] {$N$};
    \draw[F, tinta] (0.75,-0.4) -- ++(0,-0.45) node[lbl, right] {$m_p\vec g$};
    \node[lbl, font=\bfseries] at (-0.2,-1.5) {(3) pist\'on};
    \node[lbl, font=\footnotesize] at (0,-1.9) {$a_x=0$, sin rotaci\'on};
  \end{scope}
  % lectura fisica
  \node[lbl, align=left, anchor=north west, font=\small] at (10.6,2.3) {%
    \textbf{Lectura f\'isica}\\[2pt]
    $\bullet$ $N$: empuje lateral contra la camisa\\ \phantom{$\bullet$} (fricci\'on, desgaste, "piston slap")\\
    $\bullet$ $T_O$: par que entrega el cig\"ue\~nal\\
    $\bullet$ $F_{gas}$ domina a carga; la inercia\\ \phantom{$\bullet$} $m\ddot y_D\propto N^2$ domina a alta rpm};
  \node[lbl, align=center, font=\small] at (6.0,-3.2) {%
    $\sum\vec F = m(\vec a_{cm}-\vec g)$,\quad $\sum M_{cm}=I\alpha$ \;\; $\Rightarrow$\;\;
    $3+3+2=8$ ecuaciones lineales, $\mathbf{A}(\theta)\,\mathbf{x}=\mathbf{b}(\theta,\omega,F_{gas})$,\quad
    $\mathbf x=[F_O,F_C,F_D,N,T_O]$};
  \node[lbl, align=center, font=\small] at (6.0,-3.8) {%
    Verificaci\'on independiente (balance de potencia): \ $P_{eje} = F_{gas}\,(-\dot y_D) - \tfrac{d}{dt}\big(E_{cin}+E_{pot}\big)_{mecanismo}$};
\end{tikzpicture}
\end{document}
